% PDF counterparts of custom-callouts.css. Keep these accent colours in sync.
% Quarto continues to own theorem counters, titles, labels and cross-references.
\usepackage{amsthm}
\usepackage[skins,breakable]{tcolorbox}

\definecolor{courseExample}{HTML}{2D9CDB}
\definecolor{courseTheorem}{HTML}{E67E22}
\definecolor{courseDefinition}{HTML}{27AE60}
\definecolor{courseAlgorithm}{HTML}{8E44AD}
\definecolor{courseProperty}{HTML}{008080}
\definecolor{courseLemma}{HTML}{D35400}
\definecolor{courseCorollary}{HTML}{C0392B}
\definecolor{courseRemark}{HTML}{7F8C8D}
\definecolor{courseExercise}{HTML}{E74C3C}
\definecolor{courseQuestion}{HTML}{427D6B}
\definecolor{courseSolution}{HTML}{16A085}
\definecolor{courseSolutionBackground}{HTML}{F7FCFA}
\colorlet{courseEnvironmentHeading}{black}

\tcbset{course environment/.style={
  enhanced jigsaw, breakable,
  colback=white, colframe=#1, coltext=black,
  boxrule=0pt, leftrule=1.2mm, arc=1.5mm,
  left=2.5mm, right=2.5mm, top=2mm, bottom=2mm,
  before skip=0.8\baselineskip, after skip=0.8\baselineskip,
  before upper={\colorlet{courseEnvironmentHeading}{#1}}
}}

% Set titles on their own line, as in the HTML theorem blocks. Reusing
% amsthm's styles leaves Quarto's existing numbered environments intact.
\newtheoremstyle{plain}{0pt}{0pt}{\normalfont}{}
  {\bfseries\color{courseEnvironmentHeading}}{}{\newline}{}
\newtheoremstyle{definition}{0pt}{0pt}{\normalfont}{}
  {\bfseries\color{courseEnvironmentHeading}}{}{\newline}{}
\newtheoremstyle{remark}{0pt}{0pt}{\normalfont}{}
  {\bfseries\color{courseEnvironmentHeading}}{}{\newline}{}

% Install wrappers after Quarto has declared the environments. Empty hooks
% are not attached for types that do not appear in a particular render.
\newcommand{\courseBoxEnvironment}[2]{%
  \ifcsname #1\endcsname
    \tcolorboxenvironment{#1}{course environment=#2}%
  \fi
}
\AtBeginDocument{%
  % Match the current Minty HTML theme instead of Quarto's PDF defaults.
  \definecolor{quarto-callout-note-color}{HTML}{007BFF}%
  \definecolor{quarto-callout-note-color-frame}{HTML}{007BFF}%
  \definecolor{quarto-callout-tip-color}{HTML}{56CC9D}%
  \definecolor{quarto-callout-tip-color-frame}{HTML}{56CC9D}%
  \definecolor{quarto-callout-important-color}{HTML}{FF7851}%
  \definecolor{quarto-callout-important-color-frame}{HTML}{FF7851}%
  \definecolor{quarto-callout-warning-color}{HTML}{FFCE67}%
  \definecolor{quarto-callout-warning-color-frame}{HTML}{FFCE67}%
  \definecolor{quarto-callout-caution-color}{HTML}{FD7E14}%
  \definecolor{quarto-callout-caution-color-frame}{HTML}{FD7E14}%
  \courseBoxEnvironment{theorem}{courseTheorem}%
  \courseBoxEnvironment{lemma}{courseLemma}%
  \courseBoxEnvironment{corollary}{courseCorollary}%
  \courseBoxEnvironment{proposition}{courseProperty}%
  \courseBoxEnvironment{conjecture}{courseTheorem}%
  \courseBoxEnvironment{definition}{courseDefinition}%
  \courseBoxEnvironment{example}{courseExample}%
  \courseBoxEnvironment{algorithm}{courseAlgorithm}%
  \courseBoxEnvironment{exercise}{courseExercise}%
  \courseBoxEnvironment{remark}{courseRemark}%
  \courseBoxEnvironment{refremark}{courseRemark}%
  \courseBoxEnvironment{proof}{courseRemark}%
  \ifcsname solution\endcsname
    \tcolorboxenvironment{solution}{course environment=courseSolution,
      colback=courseSolutionBackground}%
  \fi
  \ifcsname refsolution\endcsname
    \tcolorboxenvironment{refsolution}{course environment=courseSolution,
      colback=courseSolutionBackground}%
  \fi
}

% Older .callout-* Divs use CSS counters in HTML. Give them the same
% per-chapter counters and accents in PDF without changing their QMD source.
\newcommand{\courseLegacyEnvironment}[3]{%
  \newcounter{courselegacy#1}[chapter]%
  \newenvironment{courselegacy#1}[1]{%
    \refstepcounter{courselegacy#1}%
    \begin{tcolorbox}[course environment=#3]%
    {\bfseries\color{#3}#2 \arabic{courselegacy#1}%
      \if\relax\detokenize{##1}\relax\else: ##1\fi\par}%
    \smallskip
  }{\end{tcolorbox}}%
}
\courseLegacyEnvironment{example}{Example}{courseExample}
\courseLegacyEnvironment{theorem}{Theorem}{courseTheorem}
\courseLegacyEnvironment{definition}{Definition}{courseDefinition}
\courseLegacyEnvironment{algorithm}{Algorithm}{courseAlgorithm}
\courseLegacyEnvironment{property}{Property}{courseProperty}
\courseLegacyEnvironment{exercise}{Exercise}{courseExercise}
\courseLegacyEnvironment{question}{Question}{courseQuestion}

% unicode-math needs a symbol alphabet for genuinely bold Greek vectors.
% Legacy \boldsymbol can otherwise select the regular math font in LuaLaTeX.
\AtBeginDocument{%
  \ifdefined\symbfit
    \renewcommand{\boldsymbol}[1]{\symbfit{#1}}%
  \fi
}

% Keep source and console output inside the one-inch text block.
\usepackage{fvextra}
\fvset{breaklines=true,breakanywhere=true,fontsize=\small}
\DefineVerbatimEnvironment{Highlighting}{Verbatim}{commandchars=\\\{\},breaklines=true,breakanywhere=true,fontsize=\small}
\RecustomVerbatimEnvironment{verbatim}{Verbatim}{breaklines=true,breakanywhere=true,fontsize=\small}
